\section{The counted finite construction}
\label{sec:extraction}

The parameters have separate roles. The profile index \(a\) determines
the limiting exponent, and fixes the denominator scale \(K\).
The type scale \(N\) reduces the relative loss in counting sector words.
The tensor power \(k\) amortizes interpolation, and the number \(j\)
of exact reuses amortizes the catalyst. We choose them in this order.

Fix \(a\ge2\) and the integer types of Section~\ref{subsec:types}.
Index the concavity and tripling relations by disjoint sets
\(\mathcal I_c\) and \(\mathcal I_t\), whose union is \(\{1,\ldots,I\}\).
At an additional positive integer scale \(N\), let the sector count of
each relation be
\begin{equation}\label{eq:sector-counts}
 M_i=\binom{Nw_i}{Nu_i}
 \quad\text{or}\quad
 M_i=\frac{(Nz_i)!}{(Nz_i/3)!^3},
\end{equation}
for a concavity or tripling relation, respectively. Define
\begin{align}
 M&=\prod_{i=1}^I M_i,&
 S&=\prod_{i=1}^I(2M_i-1)^6,& H&=M^6,\label{eq:copy-counts}\\
 D&=6\sum_{i=1}^I(M_i-1)^2,&
 u_0&=2^{2NU},&v_0&=(2^{NK}M)^2.\label{eq:dimensions}
\end{align}
The letters \(u_0,v_0\) denote matrix dimensions, not the split counts
\(u_i,d_i\).

\begin{proposition}[Counted catalytic degeneration]\label{prop:catalytic}
The specified maps give a polynomial degeneration of degree at most \(D\):
\begin{equation}\label{eq:catalytic-degeneration}
 S\Cat_N Q_{a,a}^{\otimes NK}\MM_{u_0}
 \degen H\Cat_N\MM_{v_0}.
\end{equation}
\end{proposition}
\begin{proof}
Apply Lemma~\ref{lem:determinant} with counts \((Nu_i,Nd_i)\) to
each concavity relation, and Lemma~\ref{lem:tripling} with counts
\((Nz_i/3,Nz_i/3,Nz_i/3)\) to each tripling relation.
Write this exact type map as \(\mathsf T_i\), the separator map as
\(\mathsf S_i(t)\), and the sector-specific coordinate permutations
after separation as \(\mathsf R_i\). For each \(\pi\in S_3\), the
complete relation map is
\begin{equation}\label{eq:complete-relation-map}
 \Phi_{i,\pi}(t)=\mathsf R_i^\pi\circ\mathsf S_i(t)^\pi\circ
 \left(\bigoplus_{\ell=0}^{2M_i-2}\mathsf T_i^\pi\right).
\end{equation}
Appendix~\ref{sec:composed-maps} specifies its target spaces and every
coefficient, including the positive-degree terms. In particular,
\(\mathsf R_i\) is a fixed map on the whole target space.
For every relation and every leg permutation, apply
Proposition~\ref{prop:separator} along the leg that is shared after that
permutation, with auxiliary dimension \(M_i\).
It consumes \(2M_i-1\) source copies and produces \(M_i\)
independent copies of its sector tensor, each with a dot product
of dimension \(M_i\).

Tensoring these maps over the six permutations and all relations
gives the direct-sum counts \(S\) and \(H\).
The convolution factors on the source are
\(\Cat_N Q_{a,a}^{\otimes NK}\), and those on the output are
\(\Cat_N Q_{1,2}^{\otimes NK}\), by \eqref{eq:formal-ledger}.
The neighboring-length sources contribute \(NU\) dot products of
dimension two in each of the six permutations. Their product is
\(\MM_{2^{2NU}}\).
The separator outputs contribute \(\MM_{M^2}\), while
\(Q_{1,2}^{\otimes NK}\cong\MM_{2^{2NK}}\).
Equation~\eqref{eq:matrix-products} gives the stated dimensions.

The full map is \(\bigotimes_{i,\pi}\Phi_{i,\pi}(t)\), followed by
the displayed regroupings of factors. Use one common parameter in all
separators. Lemma~\ref{lem:composition} bounds its degree by \(D\).
The initial type restrictions and final coordinate permutations are
independent of \(t\), so there are no further degree terms.
\end{proof}

The copy ratio satisfies \(1<S/H<2^{6I}\). This inequality bounds a
cost; it does not cancel either direct-sum multiplicity.

\subsection{Exactification and repeated use of one catalyst}

Define the explicit rank bound
\begin{equation}\label{eq:kappa}
 \kappa_N=\prod_{h\le b}(h+b-1)^{6Nc_{h,b}}\ge\R(\Cat_N).
\end{equation}
It comes from tensoring the decompositions
\eqref{eq:convolution-rank}, not from a rank-minimization oracle.

\begin{proposition}[An exact relative multiplication scheme]
\label{prop:relative}
For any positive integers \(k,j\), put \(E=kD+1\),
\(A=u_0^{kj}\), \(B=v_0^{kj}\), and
\begin{equation}\label{eq:relative-counts}
 r=\kappa_N^k E^j S^{kj}(2a-1)^{6NKkj},\qquad m=H^{kj}.
\end{equation}
There is a specified exact restriction \(r\MM_A\restr m\MM_B\).
\end{proposition}
\begin{proof}
Apply Lemma~\ref{lem:exactification} to the \(k\)-th power of
\eqref{eq:catalytic-degeneration}. Replace each source copy of
\(Q_{a,a}\) by its explicit rank-\((2a-1)^6\) decomposition.
Writing
\[
 \alpha=ES^k(2a-1)^{6NKk},\qquad \beta=H^k,
\]
we obtain the exact restriction
\[
 \alpha\Cat_N^{\otimes k}\MM_{u_0^k}
 \restr \beta\Cat_N^{\otimes k}\MM_{v_0^k}.
\]
This restriction may be applied \(j\) times while retaining one catalyst:
\[
 \alpha^j\Cat_N^{\otimes k}\MM_{u_0^{kj}}
 \restr\beta^j\Cat_N^{\otimes k}\MM_{v_0^{kj}}.
\]
For the induction, tensor the previous restriction with
\(\alpha\MM_{u_0^k}\), then apply the one-step restriction to the
retained catalyst in each output summand. Tensor products of matrix
multiplication tensors combine by \eqref{eq:matrix-products}.
The catalyst exponent stays \(k\), not \(kj\).

Supply its source using \(\kappa_N^k\) rank-one tensors and restrict
the output catalyst to one nonzero coefficient. A coefficient equal to
one is available by choosing a monomial in each convolution factor.
This gives \eqref{eq:relative-counts}.
\end{proof}

Exactification precedes reuse. Iterating the degeneration itself would
require a separate bound on its accumulated approximation degree.

\subsection{Batched recursion with padding}

\begin{lemma}[Relative schemes give recursive algorithms]\label{lem:batch}
Suppose an explicit restriction \(r\MM_A\restr m\MM_B\) is given,
with positive integers \(B>A\). Put
\[
 b_0=\lceil r/m\rceil,\qquad
 p_0=\frac{\log b_0}{\log(B/A)}.
\]
Then \(p_0\ge2\). The restriction yields a fixed arithmetic algorithm
for matrix multiplication with cost \(O(n^{p_0})\) if \(p_0>2\),
and \(O(n^2\log n)\) if \(p_0=2\).
Its scalar products are products of a linear form in the left input
entries and a linear form in the right input entries.
\end{lemma}
\begin{proof}
The first-leg flattenings have ranks \(rA^2\) and \(mB^2\).
Rank cannot increase under restriction, so
\(r/m\ge(B/A)^2\), which proves \(p_0\ge2\).
Compute a batch of \(m\) products of \(n\times n\) matrices. Pad each
input to side \(B\lceil n/B\rceil\) and divide it into \(B\)-by-\(B\)
blocks. The restriction reduces the batch to \(r\) matrix products of
side \(n'=A\lceil n/B\rceil\).
The scalar identity remains valid on matrix blocks: each term keeps a
left-input factor before a right-input factor, and its coefficients are
central scalars. Group the resulting products into \(b_0\) batches of
size \(m\), padding the last batch with zeros, and recurse.
All preprocessing and output combinations at a node take \(O(n^2)\)
arithmetic operations, with constants depending only on the fixed scheme.
Thus the batch cost satisfies
\[
 T(n)\le b_0T\bigl(A\lceil n/B\rceil\bigr)+C n^2
\]
above the fixed leaf threshold, for a constant \(C\).

Set \(\rho=A/B<1\) and \(c=A/(1-\rho)=AB/(B-A)\).
Then \(n'\le\rho n+A\), or \(n'-c\le\rho(n-c)\).
Choose the integer leaf threshold \(n_*=\lceil c\rceil+1\), and use classical
multiplication at leaves. At level \(\ell\), each dimension is at most
\(c+\rho^\ell(n-c)\), so there are
\(O(\log_{1/\rho}n)\) levels. Summing the work over levels gives the
two stated bounds. More explicitly, the part proportional to \(n^2\)
is bounded by a constant times
\(n^2\sum_\ell(b_0\rho^2)^\ell\); the remaining constant-size
contribution has the same or smaller order. The number of leaves is
\(O(n^{p_0})\).

Set all but one input pair of the initial batch to zero to obtain one
matrix product. The fixed batch size changes only the constant.
At each recursive call the input transformations are separately linear
in the original left and right entries. This property persists to the
classical products at the leaves, proving the last statement.
\end{proof}

In the present construction \(r/m\ge1\). Since
\(\lceil r/m\rceil\le2r/m\), Proposition~\ref{prop:relative} gives
the upper expression
\begin{equation}\label{eq:finite-exponent}
 p_0\le
 \frac{6NK\log(2a-1)+\log(S/H)+\log(kD+1)/k
       +\log\kappa_N/j+\log2/(kj)}{\log(v_0/u_0)}
\end{equation}
whenever \(v_0>u_0\).
Every copy, interpolation, catalyst, and batching charge is present.
The source \(\MM_A\) remains the recursive operation; no fast
decomposition of it has been assumed.

\subsection{Finite parameter selection}

\begin{lemma}[A quantitative type bound]\label{lem:types}
Let \(n_1+\cdots+n_s=n\), with all \(n_i>0\), and put \(q_i=n_i/n\).
Then
\[
 \frac{\exp(nH(q))}{(n+1)^{s-1}}
 \le\frac{n!}{\prod_i n_i!}\le\exp(nH(q)).
\]
\end{lemma}
\begin{proof}
The multinomial probability of the type \((n_i)\) is its coefficient
times \(\prod q_i^{n_i}=\exp(-nH(q))\), and is at most one.
It is also a mode. To see this, the factor \(n_i^h/h!\) is maximized
at \(h=n_i\) or \(h=n_i-1\); compare consecutive terms. The joint
type probability, up to a common factor, is the product of these factors.
There are at most \((n+1)^{s-1}\) types, so a mode has probability at
least their reciprocal. This proves both bounds.
\end{proof}

Put \(\Delta_N=\log(v_0/u_0)\). Equations
\eqref{eq:entropy-identity} and \eqref{eq:dimensions} imply
\begin{equation}\label{eq:type-error}
 \begin{aligned}
 2NK\log F(a,a)-2\mathcal E_N&\le\Delta_N\le2NK\log F(a,a),\\
 \mathcal E_N&=\sum_{i\in\mathcal I_c}\log(Nw_i+1)
          +2\sum_{i\in\mathcal I_t}\log(Nz_i+1).
 \end{aligned}
\end{equation}
For fixed \(a\), the error is \(O_a(\log N)\), while the main term
is linear in \(N\). As \(\log(S/H)<6I\log2\), taking \(N\) large,
then \(k\) large, and finally \(j\) large in
\eqref{eq:finite-exponent} gives the limit
\begin{equation}\label{eq:tau}
 \tau_a=\frac{3\log(2a-1)}{\log F(a,a)}\longrightarrow\frac94.
\end{equation}
For each finite \(a\), \(\tau_a>9/4\) by
Theorem~\ref{thm:least-profile}.

The limiting argument can be replaced by the following precise selection
procedure. Its comparisons involve only rational numbers raised to
integer powers; no fixed numerical precision is assumed.

\begin{theorem}[Effective finite extraction]\label{thm:effective}
There is a terminating procedure that, given a positive rational
\(\eps\), specifies a finite arithmetic algorithm for multiplying
arbitrary complex \(n\times n\) matrices in
\(O_\eps(n^{9/4+\eps})\) operations. All coefficients lie in an
explicitly specified cyclotomic field. The procedure uses only finite
integer and rational computations, coordinate permutations, tensor
products, and the maps proved above.
\end{theorem}
\begin{proof}
Write \(\theta=9/4+\eps=p/q\) in positive integers. Perform the following
searches in increasing integer order, stopping at the first successful
choice in each step.
\begin{enumerate}
\item Choose \(a\ge2\) such that
\(F(a,a)^p>(2a-1)^{3q}\). Compute the rational weights, their least
denominator \(K\), the integer types, and the catalyst counts.
\item Choose \(N\ge1\). Compute \(M_i,S,H,D,u_0,v_0,\kappa_N\) and
put \(R_0=(S/H)(2a-1)^{6NK}\). Require
\begin{equation}\label{eq:selection-N}
 (v_0/u_0)^p>R_0^q.
\end{equation}
\item With these data fixed, choose \(k\ge1\), put \(E=kD+1\), and require
\begin{equation}\label{eq:selection-k}
 (v_0/u_0)^{pk}>(E R_0^k)^q.
\end{equation}
\item Choose \(j\ge1\) such that
\begin{equation}\label{eq:selection-j}
 (v_0/u_0)^{pkj}>
 \bigl(2\kappa_N^k(E R_0^k)^j\bigr)^q.
\end{equation}
Use Proposition~\ref{prop:relative}, followed by Lemma~\ref{lem:batch}.
\end{enumerate}
Clear positive denominators to implement each test as an integer comparison.
The first search terminates by \eqref{eq:tau}. For its chosen \(a\),
\eqref{eq:type-error} and the bounded copy-ratio charge imply
\eqref{eq:selection-N} for all sufficiently large \(N\).
In particular \(v_0>u_0\), since \(R_0>1\).
In the third search the strict exponential advantage in
\eqref{eq:selection-N} dominates \(E^q=(kD+1)^q\).
Once \(k\) is fixed, the strict exponential advantage in
\eqref{eq:selection-k} dominates \((2\kappa_N^k)^q\) as \(j\)
increases. Thus every search terminates; their order keeps all previously
chosen data fixed.

The ratio of the actual counts in \eqref{eq:relative-counts} is
\(r/m=\kappa_N^k(E R_0^k)^j\). Hence
\eqref{eq:selection-j} implies
\(\lceil r/m\rceil^q<(B/A)^p\). The recursive exponent is strictly
less than \(\theta\). Any quadratic or quadratic-logarithmic term in
Lemma~\ref{lem:batch} is also \(O(n^\theta)\), because \(\theta>2\).

To specify the maps, enumerate tensor positions and sector words in
lexicographic order. Use the displayed determinant bases, interval sectors,
Fourier substitutions, evaluation decompositions, and root averages.
Separate sectors before applying their individual coordinate permutations.
Compose the maps in the order in the proofs of
Propositions~\ref{prop:catalytic} and \ref{prop:relative}.
All roots of unity have explicitly known finite orders. One may take
\[
 L_*=\lcm\bigl(E,\,2a-1,\,\{2M_i-1:1\le i\le I\},
                   \,\{h+b-1:c_{h,b}>0\}\bigr).
\]
Every coefficient lies in \(\Q(\zeta_{L_*})\); catalyst reuse and block
recursion introduce no further roots of unity.
Thus the construction specifies finite coefficient arrays as well as their
composition, even though expanding them is unnecessary for the proof.
\end{proof}

\subsection{Rational coefficients and an explicit square seed}

The field issue does not require an existence theorem for algebraic
solutions of rank equations.

\begin{lemma}[Restriction of coefficients to the rationals]
\label{lem:descent}
Let \(T\) have rational coefficients, and suppose an explicit rank-\(R\)
decomposition of \(T\) is given over a number field \(\mathbb K\) of degree
\(d\). Then an explicit rational decomposition has length at most
\(d^2R\).
\end{lemma}
\begin{proof}
Choose a rational basis \(e_1=1,e_2,\ldots,e_d\) of \(\mathbb K\), and let
\(\phi:\mathbb K\to\Q\) extract the coefficient of \(e_1\).
For each summand write its first two linear forms as
\(\sum_s e_s\alpha_s\) and \(\sum_t e_t\beta_t\), where the
\(\alpha_s,\beta_t\) have rational coefficients, and call its third
form \(\gamma\). Apply \(\phi\) coefficientwise to the tensor identity.
This summand becomes
\(\sum_{s,t}\alpha_s\beta_t\phi(e_se_t\gamma)\), a sum of at most
\(d^2\) rational rank-one tensors. The target is unchanged because its
coefficients are rational and \(\phi(1)=1\).
All field operations are effective; for the field above, use the power
basis modulo the cyclotomic polynomial.
\end{proof}

\begin{corollary}[Finite rational seeds]\label{cor:seed}
For every positive rational \(\eps\), a terminating procedure produces
an integer \(d_0\ge2\) and a rational rank decomposition of
\(\MM_{d_0}\) of length \(R_{\mathrm{seed}}<d_0^{9/4+\eps}\).
Consequently a fixed rational bilinear recursion attains an exponent
strictly below \(9/4+\eps\).
\end{corollary}
\begin{proof}
Apply Theorem~\ref{thm:effective} with \(\eps/2\).
For each integer dimension \(n\), unroll its batched recursion to
classical leaves, and set the unused initial input pairs to zero.
Lemma~\ref{lem:batch} gives a bilinear decomposition of \(\MM_n\).
Its number of leaf products is at most its arithmetic cost, hence is
\(O_\eps(n^{9/4+\eps/2})\). Its coefficient field is fixed independently
of \(n\). Apply Lemma~\ref{lem:descent}; this multiplies the length
by at most a fixed constant.

Enumerate \(n=2,3,\ldots\), generate the finite recursion and its rational
decomposition, and stop when its length \(R_n\) satisfies
\(R_n^q<n^p\), where \(p/q=9/4+\eps\).
The strict gap in exponents proves termination, regardless of the fixed
constants. This is a construction of identities by proved operations,
not a search for solutions to unspecified rank equations.
Ordinary block recursion on the resulting seed has exponent
\(\log_n R_n<9/4+\eps\); its fixed linear maps cost quadratically per node.
\end{proof}

Every rational seed identity also holds over every field of characteristic
zero. The preprocessing time of this procedure is finite for each fixed
\(\eps\). No bound uniform in \(\eps\), and no useful bit-complexity
estimate, is asserted.

For \(0<\eps\le3/4\), this prescribed seed generator returns
\(d_0>n_*>A\), where \(A,n_*\) are chosen with allowance \(\eps/2\).
At dimensions \(n\le n_*\), it produces the classical decomposition
with \(n^3\) nonzero terms, which fails the strict stopping test.
Termwise coefficient descent preserves these already rational terms.
This size constraint concerns the stated generator, not rational seeds
in general.
